% Ising quotient analysis; include after the generic CCY section.
\subsection{The Ising quotient in central-charge recursion}

It is not legitimate simply to substitute these weights into a
generic-Verma $c$-recursion and discard singular residues.  An ordinary
limit of the complete generic-Verma block is also insufficient without
a specified quotient prescription.  The problem already appears when
null states propagate around a closed cycle.  The following calculation
identifies it directly.

At $c=1/2$, the level-two null states in the modules of weights $1/16$
and $1/2$ are
\begin{equation}
 \left(L_{-2}-\frac43L_{-1}^2\right)|1/16\rangle,
 \qquad
 \left(L_{-2}-\frac34L_{-1}^2\right)|1/2\rangle.
\end{equation}
Keep the highest weights fixed and let $c$ vary.  Both displayed
states have squared norm $(c-1/2)/2$.  Indeed, in the basis
$(L_{-2}|h\rangle,L_{-1}^2|h\rangle)$ the Gram matrix is
\begin{equation}
 \begin{pmatrix}
 4h+c/2&6h\\6h&4h(2h+1)
 \end{pmatrix}.
\end{equation}
For the three-point function with highest weights
$(1/2,1/16,1/16)$, replacing either pair of entries by its displayed
null state gives $(c-1/2)/2$.  This follows either by the Virasoro Ward
identity or by identifying the $L_2$--$L_{-2}$ central contraction;
the coefficient of $c$ is $1/2$ and the value at $c=1/2$ vanishes.
The same statement holds when the middle entry is the external fermion
and the two Ramond entries are replaced by null states.

Therefore a closed cycle of level-two null states contributes with
normalized primary coefficient one.  In the vacuum term of
\eqref{eq:ising-ordered-theta-decomposition}, the first such cycle
uses all three Ramond segments.  Its leading monomial is
$\hat q_{2,1}^{2}\hat q_{2,2}^{2}q_3^2$.  In the fermion term there
are two additional cycles, with leading monomials $q_1^2q_3^2$ and
$q_1^2\hat q_{2,1}^{2}\hat q_{2,2}^{2}$.  The first cycle enters
at physical total level $4$ in the vacuum branch.  All three enter
at level $9/2$ in the fermion branch.  This is why a subtraction based
only on the Ramond character misses terms in the level-five comparison.

The subtraction involves full shifted Virasoro blocks, not just their
leading coefficients.  Write $\epsilon=c-1/2$.  A positive Virasoro
mode acting on either displayed off-critical null state gives
$O(\epsilon)$ times an ordinary state.  In a two-null three-point Ward
identity, that ordinary state is still paired with the other null
state.  The unwanted term is therefore $O(\epsilon^2)$.  Dividing by
$\epsilon$ and taking the limit gives homogeneous Ward identities at
the two shifted highest weights.  The same argument gives the shifted
descendant Gram matrix multiplied by the null norm.  This proves the
shifted-block contribution until another singular submodule enters.

Here are the formulas needed through physical total level five.  The
external weight is $1/2$ and the sewing parameters are
$(q_1,\hat q_{2,1},\hat q_{2,2},q_3)$ in every term; the parameters
are suppressed in these displays to keep the shifted weights visible.
The subscript ``irr'' specifies the irreducible block:
\begin{align}
 &F\left(0,\tfrac1{16},\tfrac1{16},\tfrac1{16};\tfrac12;\tfrac12\right)_{\rm irr}
 \nonumber\\
 &\quad=\lim_{c\to1/2}\left[
 F\left(0,\tfrac1{16},\tfrac1{16},\tfrac1{16};\tfrac12;c\right)
 -\hat q_{2,1}^2\hat q_{2,2}^2q_3^2
 F\left(0,\tfrac{33}{16},\tfrac{33}{16},\tfrac{33}{16};\tfrac12;c\right)
 \right],\label{eq:ising-vacuum-first-cycle}\\
 &F\left(\tfrac12,\tfrac1{16},\tfrac1{16},\tfrac1{16};\tfrac12;\tfrac12\right)_{\rm irr}
 \nonumber\\
 &\quad=\lim_{c\to1/2}\left[
 F\left(\tfrac12,\tfrac1{16},\tfrac1{16},\tfrac1{16};\tfrac12;c\right)
 -\hat q_{2,1}^2\hat q_{2,2}^2q_3^2
 F\left(\tfrac12,\tfrac{33}{16},\tfrac{33}{16},\tfrac{33}{16};\tfrac12;c\right)
 \right.\nonumber\\
 &\qquad\qquad\left.
 -q_1^2q_3^2
 F\left(\tfrac52,\tfrac1{16},\tfrac1{16},\tfrac{33}{16};\tfrac12;c\right)
 -q_1^2\hat q_{2,1}^2\hat q_{2,2}^2
 F\left(\tfrac52,\tfrac{33}{16},\tfrac{33}{16},\tfrac1{16};\tfrac12;c\right)
 \right].\label{eq:ising-fermion-first-cycles}
\end{align}
These equalities apply only through total physical level five,
including the factor $q_1^{1/2}$ multiplying the second one in the
auxiliary block.  In the limits, the internal and external weights
acquire distinct $O(\epsilon^3)$ perturbations; the vacuum weight is
also $O(\epsilon^3)$.  The shifts by two are applied after these
perturbations.  The complete $c$-recursive expression is evaluated
before taking its limit; individually confluent residues are never
discarded.

The degree bound needs a selection argument.  A spin null of level
four on one split segment and level-two nulls on the other segment
and the third edge would have balanced degree five.  In the vacuum
branch there is then no degree available for an NS descendant.  The
outer identity vertex between distinct shifted spin highest weights
vanishes.  In the NS fermion branch the same configuration starts at
degree $11/2$.  The NS fermion's level-three null together with a spin
level-two null also first makes a cycle at degree $11/2$; the vacuum
level-six null is already above the bound.  No additional closed null
contribution is needed in the two displayed level-five formulas.

There is a further difficulty at higher levels.  Applying the same
Ward identities with \emph{all three} entries replaced gives
\begin{align}
 &\rho\left(
 (L_{-2}-\tfrac34L_{-1}^2)|1/2\rangle,
 (L_{-2}-\tfrac43L_{-1}^2)|1/16\rangle,
 (L_{-2}-\tfrac43L_{-1}^2)|1/16\rangle
 \right)=c-\frac12.
 \label{eq:ising-triple-null}
\end{align}
Thus this vertex vanishes to first order, not second order.  If all
four internal edges carry level-two null states, the two outer
vertices give $(c-1/2)^2$, the middle vertex gives $(c-1/2)/2$, and
the four inverse Gram matrices give $16/(c-1/2)^4$.  Their contribution
therefore begins with
\begin{equation}
 \frac{8}{c-1/2}\,
 q_1^2\hat q_{2,1}^{2}\hat q_{2,2}^{2}q_3^2.
\end{equation}
After the external $q_1^{1/2}$ is restored this is physical level
$13/2$.  This fixed-weight generic-Verma limit is consequently singular
within a level-ten calculation.  Higher precision cannot remove its pole.
One must remove the interacting null submodules before claiming an
irreducible Ising block.

This does not rule out every correlated limit.  The first variation
of the triple-null vertex is twice that of a two-null vertex, and a
correlated variation can cancel both.  It must also remove the higher
primitive nulls and their intersections.  Finiteness of a full null
graph does not by itself justify an ordinary shifted primary block:
when a vertex starts at $O(\epsilon^2)$, anomalous Ward terms can
start at the same order and survive normalization.  This issue must
be resolved before extending the displayed subtraction formulas.

There is a useful exact correlated family, but it also requires further
subtractions.  Introduce a deformation parameter $x$ with $x\to-4/3$
and set
\begin{align}
 c&=13+6(x+x^{-1}),\nonumber\\
 h_L=h_R=h_3&=\frac{x+2+x^{-1}}4-\frac{x}{16},\qquad
 h_\psi=-\frac12-\frac{3x}{4}.
 \label{eq:ising-exact-family}
\end{align}
The external field is $(2,1)$ throughout this family.  Its degenerate
fusion maps the Ramond momentum to its negative and therefore preserves
the Ramond weight.  In the NS fermion branch one may also set
$h_1=h_\psi$ and first quotient its generic level-two null submodule.
Since that NS edge alone is not a cycle, its independent preliminary
quotient does not leave a closed null contribution.  The analogous
preliminary operation removes $L_{-1}|0\rangle$ in the vacuum branch.

Nevertheless, at $c=1/2$ the NS fermion module develops its second
primitive null at level three.  The explicit states relevant to this
remaining problem are
\begin{align}
 &\left(\tfrac34 L_{-3}-3L_{-2}L_{-1}+L_{-1}^3\right)|1/2\rangle,
 \nonumber\\
 &\left(-\tfrac14 L_{-4}+\tfrac{11}{6}L_{-3}L_{-1}
 +\tfrac{49}{144}L_{-2}^2-\tfrac{25}{6}L_{-2}L_{-1}^2
 +L_{-1}^4\right)|1/16\rangle.
\end{align}
They are the level-three NS null and the level-four Ramond null,
respectively; the level-two Ramond null was displayed above.  Along
\eqref{eq:ising-exact-family}, the first derivatives of their squared
norms with respect to $c$ are $75/8$ and $1925/648$.  The corresponding
level-two Ramond norm derivative is $4/9$.  The three-point Ward
identities give first derivatives $5/4$ and $275/48$ when the NS
level-three null and, respectively, the Ramond level-two or level-four
null are inserted at an outer vertex.  Exchanging Ramond slots can
reverse a sign, but both original vertices have the same ordering.
The closed null cycle along the NS edge and the third edge consequently
has leading coefficients
\begin{equation}
 \frac{(5/4)^2}{(75/8)(4/9)}=\frac38,
 \qquad
 \frac{(275/48)^2}{(75/8)(1925/648)}=\frac{33}{28}.
\end{equation}
Their monomials are $q_1^3q_3^2$ and $q_1^3q_3^4$ before the
NS factor $q_1^{1/2}$ is restored.  They therefore survive at physical
levels $11/2$ and $15/2$.  The preliminary generic NS quotient does not
change these coefficients: adding a descendant of its always-null
level-two vector changes neither three-point form on the exact fusion
family.  Thus \eqref{eq:ising-exact-family} alone is also insufficient
for level ten, even though it improves the first interacting null
singularity.  This is a separate, finite counterexample to that
prescription, not a claim that every correlated prescription fails.

The implementation exposes the exact parity assembly and the
first-cycle subtraction restricted to total level five.  Its production
API enforces this bound.  Both authorized integrated comparisons have
passed for that prescription; the results appear in the main
numerical-status section.  An uncorrected generic CCY limit remains
available only as an explicitly unproved candidate.  None of these is
declared a completed level-ten Ising algorithm.
The primitive nulls at levels $3$ in the fermion module, $4$ in the
spin module, and $6$ in the vacuum module, together with intersections
of null submodules and the interaction in
\eqref{eq:ising-triple-null}, still require a complete quotient
prescription.  The ordinary vacuum null $L_{-1}|0\rangle$ must also
be removed before the rational-central-charge limit.

The exact symbolic derivatives used here, including those of all
primitive three-null vertices relevant through level ten, are derived
in \texttt{ising\_tangent\_derivation.py} and documented in
\texttt{ISING\_TANGENT\_ANALYSIS.md}, both under
\texttt{Code/theta\_fermion\_ccy}.  They are
algebraic Ward calculations, not additional numerical block comparisons.

For related warnings about path-dependent rational limits, including
the distinction between Verma and irreducible characters, see
N.~Javerzat, R.~Santachiara, and O.~Foda,
\emph{Notes on the solutions of Zamolodchikov-type recursion relations
in Virasoro minimal models}, arXiv:1806.02790, Secs.~3 and~5.
Their sphere conjecture and torus character prescription do not by
themselves provide the missing punctured genus-two quotient.
